\documentclass[11pt]{article}
\usepackage[margin=1.05in]{geometry}
\usepackage{amsmath,amssymb}
\usepackage{booktabs}
\usepackage{enumitem}
\usepackage[hidelinks]{hyperref}

\title{Stage 2: Long-Horizon MimicGen Tasks\\
\large Why These Tasks, the Two-Simulator Problem, the Training Setup, and Current Status}
\author{Aryan Goyal}
\date{July 20, 2026 (updated July 22)}

\begin{document}
\maketitle

\begin{abstract}
\noindent
Stage 1 ended with a precise open question. On short, contact-dominated tasks the
predictor-driven chunk switch ties the best fixed chunk length but does not beat
it. The theory says the win should appear on long tasks that alternate between
free-space transport and contact-rich stages, because there no single fixed
chunk length can be right everywhere. Stage 2 tests exactly that. It trains
Diffusion Policy on four long-horizon MimicGen tasks (three\_piece\_assembly,
nut\_assembly, kitchen, coffee\_preparation, all D0, 1{,}000 demonstrations
each, mean episode lengths 337 to 689 steps) and will then run the same
fixed-$k$ versus predictor-switch comparison as Stage 1c. This document records
the design decisions, the two-simulator version problem and how it was solved,
the exact conversion and training pipeline, the status log (including the
July 21 scratch-disk event that killed the first training runs at 80 to 86
percent, the relaunch fault that then deleted the surviving checkpoints, and
the hardened restart that now carries the pipeline), and the open decisions
for the evaluation phase.
\end{abstract}

\tableofcontents
\newpage

%=====================================================================
\section{Why Stage 2 exists}
\label{sec:why}

Stage 1c's pooled result is the starting point. On tool\_hang the predictor
switch reached 0.800 over three seeds, which statistically ties both the best
fixed chunk length (0.807) and the best contact-oracle cell (0.807), while
avoiding the 0.60 to 0.67 failure of per-step replanning in contact. That is a
tie, not a win, and the Stage 1a analysis explains why a tie may be the best
possible outcome on such tasks: square and tool\_hang rollouts spend 65 to 88
percent of their steps in contact, so a single well-chosen fixed $k$ is nearly
optimal and an adaptive switch has almost nothing to modulate.

The hypothesis makes a sharper prediction for tasks where no single $k$ can be
right. A long task that alternates several times between free-space transport
(stable, wants long chunks) and contact-rich manipulation (unstable, wants
short chunks) forces every fixed $k$ into a compromise: short $k$ pays the
per-step replanning cost through the long stable stretches, long $k$ carries
open-loop error through the contact stages. An adaptive switch pays neither
cost. Stage 2 is the fair test of that prediction.

Two properties matter when picking the tasks:

\begin{enumerate}[leftmargin=2em]
  \item \textbf{Length and structure.} Episodes should be several hundred steps
        with multiple distinct manipulation stages separated by transport, so
        the stable and unstable segments both occupy meaningful fractions of
        the episode.
  \item \textbf{Headroom.} Base-policy success should sit well below 1.0, so an
        execution-rule improvement is measurable. Stage 1a's square result
        showed what happens without headroom: everything lands in a flat band.
\end{enumerate}

%=====================================================================
\section{The four tasks}
\label{sec:tasks}

MimicGen (Mandlekar et al., 2023) provides 1{,}000-demonstration datasets for
multi-stage robosuite tasks, generated automatically from small sets of human
source demonstrations. We use the D0 variant of four core tasks, which is the
default object-initialization distribution. The datasets ship as state-only
trajectories (simulator states plus actions), which is why the conversion step
in Section \ref{sec:convert} exists.

\begin{center}
\small
\begin{tabular}{@{}lrrrr@{}}
\toprule
Task & Demos & Transitions & Mean length & Length range \\
\midrule
three\_piece\_assembly\_d0 & 1{,}000 & 336{,}695 & 337 & 308 to 375 \\
nut\_assembly\_d0          & 1{,}000 & 358{,}907 & 359 & 293 to 435 \\
kitchen\_d0                & 1{,}000 & 616{,}751 & 617 & 585 to 654 \\
coffee\_preparation\_d0    & 1{,}000 & 689{,}273 & 689 & 591 to 760 \\
\bottomrule
\end{tabular}
\end{center}

In plain terms: three\_piece\_assembly inserts two pieces onto a base in
sequence; nut\_assembly places both a square and a round nut onto their pegs;
kitchen runs a multi-stage cooking sequence involving a stove, a pot, and
bread; coffee\_preparation loads a mug and prepares the coffee machine across
several steps. Each task chains grasp, transport, and precise-placement stages,
which is exactly the alternating structure Section \ref{sec:why} asks for. For
scale, the Stage 1 tasks were one-stage or two-stage tasks with 200
demonstrations each; these are two to five times longer per episode, have five
times the demonstrations, and published imitation-learning success rates on
them sit well below ceiling.

%=====================================================================
\section{The two-simulator problem}
\label{sec:venv}

Stage 2 has one structural complication that shaped everything else, so it gets
its own section.

MimicGen environments require robosuite 1.4 (they are registered by the
\texttt{mimicgen} package, which pins the older simulator API). Our Stage 1
policy-training stack lives on the current robomimic main branch, which is the
only branch with the Diffusion Policy implementation, and it runs against
robosuite 1.5. The two requirements cannot coexist in one Python environment.
Additionally, the robomimic v0.3 release that matches the MimicGen ecosystem
has no Diffusion Policy at all (its template set stops at BC, BCQ, CQL, HBC,
IQL, IRIS, TD3-BC).

The resolution is two virtual environments with a strict division of labor:

\begin{itemize}[leftmargin=2em]
  \item \textbf{The mg venv} (robosuite 1.4.1, mujoco 2.3.2, robomimic v0.3,
        mimicgen): everything that must instantiate a MimicGen environment.
        Today that is dataset conversion; later it is policy evaluation.
  \item \textbf{The lh venv} (robosuite 1.5, robomimic main): everything that
        touches Diffusion Policy. Today that is training; it can never create a
        MimicGen environment.
\end{itemize}

The immediate casualty is in-training rollout evaluation. Stage 1 trainings
rolled out the policy every 50 epochs to track success; Stage 2 trainings
cannot, because the training process (lh venv) cannot build the environment.
Training therefore runs with rollouts disabled, and success is measured post
hoc from saved checkpoints through the mg venv (Section \ref{sec:eval}). This
was validated up front with a one-epoch smoke test of the full training
pipeline on a converted MimicGen dataset, which completed cleanly.

%=====================================================================
\section{Step 1: dataset conversion}
\label{sec:convert}

The MimicGen hdf5 files store simulator states, not observations. The robomimic
v0.3 script \texttt{dataset\_states\_to\_obs.py} replays each demonstration
inside the environment and renders the observations the policy will train on.
Per task we render:

\begin{itemize}[leftmargin=2em]
  \item Two cameras, \texttt{agentview} and \texttt{robot0\_eye\_in\_hand}, at
        84 by 84 pixels. This matches the robomimic image recipe the Stage 1
        policies used.
  \item The standard low-dimensional states (end-effector pose, gripper).
  \item \texttt{--exclude-next-obs}, which roughly halves the file size;
        behavior cloning never reads next-observations.
  \item \texttt{--done\_mode 2}, the convention used by the released robomimic
        image datasets.
\end{itemize}

Two mechanical details worth recording. First, the script only works on
MimicGen data if \texttt{import mimicgen} runs before the environment is
created, because that import registers the environment classes; the script has
no hook for this, so we invoke it through a small \texttt{runpy} wrapper that
performs the import, sets \texttt{sys.argv}, and then executes the script as
\texttt{\_\_main\_\_}. Second, rendering is headless EGL, pinned per GPU with
\texttt{MUJOCO\_GL=egl} and \texttt{MUJOCO\_EGL\_DEVICE\_ID}; EGL processes do
not appear in \texttt{nvidia-smi}'s compute-apps listing, so the reliable
liveness check is output-file growth, not GPU process tables.

Each of the four conversions ran on its own GPU (0 through 3), writing
\texttt{<task>\_image.hdf5} next to the source file and logging to
\texttt{<task>\_img\_gen.log}. Measured durations and output sizes:

\begin{center}
\small
\begin{tabular}{@{}lrr@{}}
\toprule
Task & Conversion time & Output size \\
\midrule
nut\_assembly\_d0          & 78 min  & 14.9 GB \\
kitchen\_d0                & 92 min  & 25.7 GB \\
three\_piece\_assembly\_d0 & 104 min & 14.1 GB \\
coffee\_preparation\_d0    & 119 min & 29.0 GB \\
\bottomrule
\end{tabular}
\end{center}

Sizes track transition count times the uncompressed two-camera image payload
(about 42 KB per step); the full set is 84 GB against 2.4 TB free on the
scratch disk.

%=====================================================================
\section{Step 2: policy training}
\label{sec:train}

Each GPU chain launches its Diffusion Policy training automatically the moment
its conversion finishes, so no idle time sits between the steps. The training
configuration is generated by \texttt{impl/configs/gen\_mg\_configs.py} from
the robomimic \texttt{diffusion\_policy.json} template, and it deliberately
changes as little as possible from the Stage 1 Panda trainings:

\begin{itemize}[leftmargin=2em]
  \item \textbf{Identical budget.} The template sets
        \texttt{epoch\_every\_n\_steps} to 100, so 2{,}000 epochs is 200{,}000
        gradient steps regardless of dataset size. This is the same budget the
        Panda policies got, which keeps later comparisons about the tasks, not
        the training length.
  \item \textbf{Same observation recipe.} RGB \texttt{agentview} plus
        \texttt{robot0\_eye\_in\_hand} through a 76 by 76 CropRandomizer;
        low-dimensional end-effector position, quaternion, and gripper state.
  \item \textbf{Rollouts and validation disabled.} Rollouts for the reason in
        Section \ref{sec:venv}. Validation because the MimicGen hdf5s ship
        without train/valid mask groups, so the filter keys are set to none and
        all 1{,}000 demonstrations train.
  \item \textbf{Checkpoints every 50 epochs}, giving 40 checkpoints per task
        for the post-hoc evaluation to scan.
  \item \textbf{Dataset cached in RAM} (\texttt{hdf5\_cache\_mode: all}); the
        converted files fit comfortably in this machine's memory.
\end{itemize}

The measured pace is much faster than the Stage 1 trainings, and the reason is
instructive. The Stage 1 wall-clock (days per task) was dominated by
in-training rollout evaluation; with rollouts disabled and the dataset cached
in RAM, an epoch here is just 100 gradient steps at about 8 steps per second
plus bookkeeping, or 30 to 45 wall seconds. At that pace 2{,}000 epochs is 17
to 24 hours per task, and since all four run concurrently, the whole training
phase fits in a day. The venv constraint that forced rollouts off, in other
words, also collapsed the schedule.

%=====================================================================
\section{Step 3, pending: post-hoc evaluation}
\label{sec:eval}

Measuring success requires rolling out lh-venv-trained checkpoints inside
mg-venv environments, and no single process can do both. Two candidate designs,
not yet decided:

\begin{enumerate}[leftmargin=2em]
  \item \textbf{Policy server.} The policy loads in an lh-venv process; an
        mg-venv process owns the environment and exchanges observations and
        actions with the policy process over a local socket or pipe. This is
        the same pattern Stage 1c used to run the predictor (which needs a
        newer transformers stack) next to the rollout process, so the
        infrastructure precedent exists in this codebase.
  \item \textbf{Backport.} Install a Diffusion-Policy-capable robomimic into
        the mg venv if its torch and robosuite pins allow it. Cleaner if it
        works, but the dependency solve is unverified and the policy-server
        pattern is known-good.
\end{enumerate}

The bridge is the immediate next work item: with the measured training pace
(Section \ref{sec:train}) the full checkpoint sets exist by the morning of
July 21, so the bridge, not the trainings, is the critical path. Its first use
is checkpoint selection (success rate at a subsample of the 40 checkpoints),
mirroring how the Stage 1 policies were selected.

%=====================================================================
\section{Step 4, pending: the Stage 2 experiment}
\label{sec:experiment}

Once each task has a selected checkpoint, the experiment repeats the Stage 1c
protocol on the long-horizon tasks: fixed $k \in \{1, 4, 8, 16\}$ and the
predictor-driven $(16, 4)$ switch, same policy weights across all cells, 50
episodes per seed, horizons set per task from the demonstration lengths. The
prediction being tested: on these tasks the adaptive switch beats every fixed
$k$, because the alternating stable/unstable structure makes every fixed choice
pay one of the two costs from Section \ref{sec:why}.

One decision is still open and it is the main scientific choice left in Stage
2: what the predictor head is at evaluation time.

\begin{enumerate}[leftmargin=2em]
  \item \textbf{Zero-shot.} Use the frozen Stage 1 head as-is. The DINOv2
        ablation showed the head reads scene configuration, and the MimicGen
        scenes are robosuite scenes visually adjacent to the training tasks, so
        transfer is plausible. This is the cheap default and the more
        impressive result if it works.
  \item \textbf{Relabel and fine-tune.} Run the FTLE labeler over the MimicGen
        demonstrations (through the mg venv) and train a head that has seen
        these tasks. This is days of labeling compute and requires adapting the
        labeler to the mg environments, but it separates ``the switch idea
        works on long-horizon tasks'' from ``the head transfers zero-shot,''
        which are different claims.
  \end{enumerate}

The honest sequencing is to run zero-shot first, since it needs nothing new,
and to decide on relabeling based on what the zero-shot head's predictions look
like on MimicGen frames (sane $\hat\lambda$ distributions and plausible
switching would justify skipping the relabel; degenerate always-long or
always-short output would force it).

That check has now run (July 20, 320 stamps per task sampled from 40
demonstrations each, agentview clips plus proprioception exactly as in
training):

\begin{center}
\small
\begin{tabular}{@{}lcccc@{}}
\toprule
 & mean $\hat\lambda$ & p5 to p95 & frac $> \delta$ & corr with episode time \\
\midrule
three\_piece\_assembly & 0.022 & $-0.008$ to $0.046$ & 0.19 & $-0.04$ \\
nut\_assembly          & 0.017 & $-0.016$ to $0.050$ & 0.22 & $0.00$ \\
kitchen                & 0.019 & $-0.009$ to $0.038$ & 0.11 & $0.16$ \\
coffee\_preparation    & 0.035 & $0.002$ to $0.058$  & 0.55 & $0.26$ \\
\bottomrule
\end{tabular}
\end{center}

The distributions are varied, span the training threshold $\delta = 0.0355$ on
every task, and would drive real switching (11 to 55 percent of stamps land
above $\delta$). They are also ordered plausibly: coffee\_preparation, the most
contact-heavy manipulation sequence, reads most unstable, while kitchen, which
is dominated by transport between fixtures, reads most stable. The decision is
therefore zero-shot first. Two limits of this check are worth keeping in mind:
it ran on demonstration frames, not on frames under policy control, and it
verifies non-degeneracy, not accuracy; no ground-truth labels exist on these
tasks, so the rollout comparison itself remains the real test. Relabeling
stays available as the follow-up if the zero-shot switch underperforms.

%=====================================================================
\section{Status log, July 20 to 22, 2026}
\label{sec:status}

\begin{itemize}[leftmargin=2em]
  \item \textbf{09:05 UTC.} All four conversions launched, one per GPU, after
        the pipeline smoke test passed.
  \item \textbf{10:23 to 11:04 UTC.} Conversions completed in launch order of
        their durations (nut\_assembly 10:23, kitchen 10:37,
        three\_piece\_assembly 10:49, coffee\_preparation 11:05), and each
        Diffusion Policy training auto-launched within a second of its
        conversion finishing. No failures.
  \item \textbf{11:25 UTC.} All four trainings verified stepping: epoch
        progress at about 8 gradient steps per second on every GPU, roughly
        4 GB of GPU memory each, 21 to 95 epochs completed, and the first
        periodic checkpoint (epoch 50) already on disk for nut\_assembly.
  \item \textbf{Projection.} At the measured 30 to 45 seconds per epoch, the
        four trainings complete the morning of July 21 UTC, staggered in
        launch order.
  \item \textbf{12:15 UTC.} The evaluation bridge (Section \ref{sec:eval})
        built and smoke-tested end to end on a partially trained
        nut\_assembly checkpoint: full protocol round trip at about 7.7
        environment steps per second, which prices a 50-episode cell at
        roughly 70 minutes on one GPU. Three cross-venv format mismatches
        were found and fixed in the process (stdout noise in the protocol
        stream, v0.3-processed images versus main-branch raw-image
        expectations, and the checkpoint's frame stacking, which the client
        now reproduces).
  \item \textbf{12:50 UTC.} The zero-shot distribution check (Section
        \ref{sec:experiment}) ran and passed; the head is not degenerate on
        MimicGen frames and the zero-shot-first decision is taken.
  \item \textbf{Remaining before the experiment.} Checkpoint selection once
        trainings finish (morning of July 21), then the fixed-$k$ sweep plus
        predictor cells.
\end{itemize}

%=====================================================================
\section{July 21: the disk event and the restart}
\label{sec:incident}

Overnight between July 20 and 21 the scratch disk was largely wiped: the
source and converted datasets, the three Python environments, the cloned
repositories, and the extracted feature files were deleted. The four training
processes died at 05:01 UTC, at epochs 1,721 (nut\_assembly), 1,653 (kitchen),
1,637 (three\_piece\_assembly), and 1,596 (coffee\_preparation) of 2,000. They
had survived the deletion itself for several hours (their code and RAM-cached
data were held as open file handles) and then went down without a Python
traceback, consistent with the disk being reprovisioned under them.

What survived, and why the damage is small:

\begin{itemize}[leftmargin=2em]
  \item \textbf{All checkpoints and results} live on persistent storage and
        were untouched. Checkpoints saved every 50 epochs, so the surviving
        grids end at epoch 1,700 / 1,650 / 1,600 / 1,550 per task. All four
        are far past the epoch 950 to 1,200 region where Stage 1's best
        checkpoints lived, so these are almost certainly at-plateau policies.
  \item \textbf{The labels} (31 npz files, days of GPU compute) survived on
        scratch and were immediately copied to persistent storage, along with
        the lift and can rollout datasets. The square and tool\_hang rollout
        datasets were lost; they are only needed for future feature work, not
        for Stage 2.
  \item \textbf{Code and documents} live on persistent storage; nothing was
        lost.
\end{itemize}

The rebuild (venvs, repositories, dataset re-downloads, image reconversion)
hit two reproducible faults worth recording. First, the fresh mg venv resolved
a torchvision wheel that did not match its CUDA torch build; the fix is to
install torch and torchvision together from the same CUDA index. Second, the
fresh robomimic v0.3 clone restored an unconditional \texttt{import mujoco\_py}
in its robosuite environment wrapper; the original clone had carried a local
try/except patch that was wiped with it. The patch is now written into the
setup script so the next rebuild applies it automatically.

The restart decision, made at 12:21 UTC: run both tracks in parallel rather
than serialize them.

\begin{itemize}[leftmargin=2em]
  \item \textbf{GPUs 0 to 3}: reconvert the datasets (about two hours), then
        fresh 2,000-epoch training runs, completing around midday July 22.
        These reruns are completion insurance; nothing on the critical path
        waits for them.
  \item \textbf{GPUs 4 to 7}: checkpoint selection over the surviving
        checkpoint grids (epochs 600, 1,000, 1,400, and the last saved
        checkpoint per task), starting per task as its reconversion lands and
        finishing the evening of July 21. The Stage 2 experiment launches on
        the selected checkpoints, so the experiment timeline moves by less
        than a day despite the incident.
\end{itemize}

\subsection{July 21, afternoon: the relaunch destroyed the surviving
checkpoints}
\label{sec:incident2}

The parallel-tracks plan failed in its first hour, through the relaunch
script rather than the disk. robomimic's training entry point prompts for
confirmation when its experiment directory already exists, and answering yes
makes it delete that directory, checkpoints included. The relaunch command
piped an automatic yes into that prompt. The result: the surviving checkpoint
grids for all four tasks (epochs up to 1,700 / 1,650 / 1,600 / 1,550) were
deleted at relaunch, and the checkpoint selection pass that then ran found no
checkpoint files and completed vacuously. The loss is permanent, because the
same-day document backup had deliberately excluded the training outputs for
size.

The fresh 2,000-epoch reruns, launched 13:39 to 14:19 UTC on the same
datasets and configs, went from completion insurance to the pipeline itself.
Three process changes came out of the failure. First, the launch script now
refuses to start a training when the experiment directory exists, instead of
answering the overwrite prompt; a human has to delete the directory
deliberately. Second, checkpoint selection was rewritten to wait for each
training to report success and then copy that task's full checkpoint grid to
the OS disk before evaluating anything, so no future fault can destroy the
only copy. Third, the vacuous selection marker was cleared so the automatic
shutdown watchdog (armed at the user's request) cannot fire on it.

Net cost of the second incident: about one day. As of 06:30 UTC July 22 the
four fresh runs are healthy at epochs 1,425 to 1,643 of 2,000, on pace to
finish between roughly 10:00 and 13:00 UTC. Selection (epochs 600, 1,000,
1,400, 1,700, 2,000 per task, 25 episodes each) starts per task as its
training lands, and the switch-versus-fixed-$k$ experiment follows on the
selected checkpoints.

%=====================================================================
\section{Risks and open items}
\label{sec:risks}

\begin{itemize}[leftmargin=2em]
  \item \textbf{Zero-shot head accuracy is unverified.} The planned
        degeneracy check has run and passed (Section \ref{sec:experiment}),
        so the switch will actually switch; whether it switches at the right
        moments is untested, because no ground-truth labels exist on these
        tasks. The distribution check also used demonstration frames, and
        frames under policy control may differ. Fallback if the switch
        underperforms: relabel and fine-tune.
  \item \textbf{Base-policy success is unknown until the bridge exists.} If a
        task trains to very low success, the switch comparison on it loses
        power. Four tasks were chosen partly so that one weak task does not
        sink the stage.
  \item \textbf{No in-training signal.} With rollouts and validation disabled,
        training health is only visible through the loss curve until the bridge
        runs. The losses are being watched; divergence would be caught within a
        checkpoint interval.
  \item \textbf{Horizon choices for evaluation.} Rollout horizons must be long
        enough that slow-but-successful episodes are not truncated; they will
        be set from the demonstration length maxima with margin, per task.
  \item \textbf{Deferred relative to the original plan.} The original Stage 2
        (ACT on ALOHA, execution-rule ablations) is deferred, not dropped; the
        long-horizon test was promoted because Stage 1c identified it as the
        decisive missing experiment. FurnitureSim remains parked on a manual
        download. LIBERO-Long policy training remains an open option for the
        idle GPUs.
\end{itemize}

%=====================================================================
\section{Summary}

Stage 1 showed the switch matches the best fixed chunk length on short
contact-dominated tasks and explained why it cannot beat it there. Stage 2 asks
the question the theory actually cares about, on four MimicGen tasks whose
episodes are long, multi-stage, and alternate between transport and contact.
The two-simulator conflict between MimicGen (robosuite 1.4, no Diffusion
Policy) and our training stack (robosuite 1.5) is solved by a strict
conversion-in-mg, training-in-lh split, at the cost of moving all evaluation
post hoc. The first training runs were killed at 80 to 86 percent by the
July 21 scratch-disk event (Section \ref{sec:incident}), and the surviving
checkpoint grids were then deleted by a fault in the relaunch script itself
(Section \ref{sec:incident2}), so the fresh full-length reruns launched on
July 21 now carry the whole pipeline. They are healthy and finish around
midday July 22 UTC. The evaluation bridge is built and smoke-tested, the
zero-shot decision for the predictor head is taken, and the remaining
sequence is: per-task checkpoint backup and selection as each training
finishes (July 22), then the fixed-$k$ sweep plus predictor cells (July 22
to 23), then automatic shutdown.

\end{document}
